% <!-- PRIMARY-SOURCE-EXEMPT: reason=paper with existing bibliography (wildberger2005, barning1963, dale2026fib, wall1960, carmichael1910) -->
\documentclass[11pt]{article}

\usepackage{amsmath,amssymb,amsthm}
\usepackage{booktabs}
\usepackage[margin=2.5cm]{geometry}
\usepackage{hyperref}
\usepackage{enumitem}

\newtheorem{definition}{Definition}[section]
\newtheorem{proposition}[definition]{Proposition}
\newtheorem{theorem}[definition]{Theorem}
\newtheorem{corollary}[definition]{Corollary}
\newtheorem{remark}[definition]{Remark}
\newtheorem{example}[definition]{Example}

\newcommand{\Zm}{\mathbb{Z}/m\mathbb{Z}}
\newcommand{\Zphi}{\mathbb{Z}[\varphi]}
\newcommand{\QA}{\mathrm{QA}}

\begin{document}

\title{Modular Fibonacci Orbit Classification:\\
Convergent Structure Across Graph, Planetary, Neural, and Climate Domains}
\author{Will Dale\\
\small Independent Researcher}
\date{}
\maketitle

\begin{abstract}
We report that a single discrete algebraic object --- the orbit decomposition of the
Fibonacci shift $T(b,e)=(e,((b+e-1)\bmod m)+1)$ on the state space
$\{1,\dots,m\}^2$ at modulus $m=24$ --- produces statistically significant structure
in four independent empirical domains: graph network topology, planetary mean-motion
resonances, EEG seizure prediction, and climate regime classification.
In the graph-theoretic case we prove an exact theorem (Anchor Geodesic Separation):
for any finite unweighted tree, the sign product $\Delta b \cdot \Delta e$ of
BFS-distance increments separates on-path from off-path edges without error.
In the three empirical domains the same orbit classification --- applied as an
observer projection with no parameter fitting in the algebraic layer --- yields
$p < 10^{-6}$ (planetary, out-of-sample $p < 10^{-45}$), $\Delta R^2 = +0.21$
(EEG, $p < 10^{-32}$), and $\chi^2 = 95$ ($p < 10^{-6}$, climate).
The convergence is not post-hoc: the same modulus $m=24$ arises independently
as the Pisano period $\pi(9)$ and as a divisibility constraint on primitive
Pythagorean triples; it is not tuned to any of the four domains.
We discuss the observer-projection firewall that prevents circular reasoning,
the tree-specificity boundary of the formal theorem, and the alternative
hypothesis that the orbit classifier is an efficient nonlinear feature extractor
whose cross-domain success is coincidental.
\end{abstract}

% ====================================================================
\section{Introduction}\label{sec:intro}

Cross-domain applicability of a mathematical structure is evidence of a kind
different from within-domain fit.  Within a single domain, tuning --- conscious
or unconscious --- can produce spurious results.  When the same unmodified
algebraic object predicts anomalies in graph networks, planetary mechanics,
neural signals, and climate, the tuning hypothesis requires the object to have
been independently over-fitted to four domains, which is a strong prior to overcome.

The Quantum Arithmetic (QA) system is a discrete dynamical system on the pair space
$S_m = \{1,\dots,m\}^2$ generated by the Fibonacci shift
$T(b,e) = (e,\, ((b+e-1) \bmod m) + 1)$.
This operates in $\mathbb{Z}[\varphi]/m\mathbb{Z}[\varphi]$ where $\varphi = (1+\sqrt5)/2$,
and $T$ corresponds to multiplication by $\varphi^2$.
For $m=9$ the orbit decomposition produces three types:
\begin{itemize}
  \item \textbf{Cosmos}: the 72-element main cycle (period 24 within mod-9 arithmetic).
  \item \textbf{Satellite}: an 8-element sub-cycle with distinct geometric structure.
  \item \textbf{Singularity}: the fixed point $(9,9)$ under $T$.
\end{itemize}
The canonical applied modulus $m=24$ is the Pisano period $\pi(9)$:
the Fibonacci sequence returns to $(1,1)$ modulo $9$ every 24 steps.
This modulus arises independently from the divisibility constraint
$24 \mid 2CF$ for every primitive Pythagorean triple $(C,F,G)$.
Neither origin involves the empirical domains studied here.

The methodological backbone is what we call the \emph{observer-projection firewall}:
continuous or domain-specific measurements enter the algebraic layer exactly once
(as a discretisation map), the QA dynamics produce a discrete orbit label,
and that label is then compared against an outcome.
The continuous measurement never serves as a causal input to QA dynamics.
This prevents the circular reasoning that would arise if, for example,
a seizure probability were used to set the QA state which then ``predicts''
the seizure probability.

In Section~\ref{sec:ags} we state and prove the Anchor Geodesic Separation (AGS) theorem,
which gives the graph-theoretic result an exact algebraic proof.
Sections~\ref{sec:planetary}--\ref{sec:climate} report the three empirical domains.
Section~\ref{sec:synthesis} discusses cross-domain synthesis and alternative explanations.

% ====================================================================
\section{The Algebraic Structure}\label{sec:algebra}

\begin{definition}[QA state and orbit]
Fix modulus $m \geq 2$.  The \emph{state space} is $S_m = \{1,\dots,m\}^2$.
The \emph{Fibonacci shift} is $T(b,e) = (e,\, d)$ where $d = ((b+e-1)\bmod m)+1$.
For any state $(b,e)$, the orbit $\mathcal{O}(b,e) = \{T^k(b,e) : k \geq 0\}$
falls into one of three types at $m=9$: Cosmos (period 24), Satellite (period 8),
or Singularity (the fixed point $(9,9)$).
\end{definition}

\begin{definition}[Observer projection]
A function $\phi : X \to S_m$ from a domain-specific space $X$ (e.g.\ spectral
measurements, orbital periods, BFS distances) is an \emph{observer projection} if
it maps domain data into QA state coordinates without using QA dynamics to define it.
The QA output for $x \in X$ is the orbit type of $\phi(x)$.
\end{definition}

The key property exploited across all four domains is that the orbit type is a
\emph{discrete invariant} of the Fibonacci dynamics: once a pair $(b,e)$ is
classified as Cosmos, Satellite, or Singularity, no continuous deformation of
the input (short of crossing a discretisation boundary) changes the classification.
This makes the orbit label structurally robust in a way that continuous features are not.

% ====================================================================
\section{Anchor Geodesic Separation}\label{sec:ags}

The graph-theoretic application assigns BFS distances to two anchor nodes as the
$(b,e)$ coordinates of each graph node.  This is the cleanest instance of the
observer-projection pattern: the BFS distances are external measurements of graph
structure, the $\Delta b \cdot \Delta e$ sign product is a purely integer computation,
and the output (on-path vs.\ off-path) is a discrete binary label.

\begin{definition}[BFS coordinates]
Let $T = (V,E)$ be a finite unweighted tree with anchor nodes $L, R \in V$.
For $v \in V$ define $b(v) = d_T(v,L)$ and $e(v) = d_T(v,R)$.
For edge $(v,u)$ write $\Delta b = b(u)-b(v)$ and $\Delta e = e(u)-e(v)$.
\end{definition}

\begin{theorem}[Anchor Geodesic Separation]\label{thm:ags}
For every edge $(v,u) \in E$: (i)~$\Delta b, \Delta e \in \{-1,+1\}$;
(ii)~$\Delta b \cdot \Delta e = -1$ iff the edge lies on the unique $L$--$R$ path;
(iii)~$\Delta b \cdot \Delta e = +1$ iff the edge does not lie on the $L$--$R$ path.
\end{theorem}

\begin{proof}
(i) Any tree edge lies on exactly one geodesic from $L$; distance changes by exactly one.
(ii--iii) On the $L$--$R$ path, each step moves toward one anchor and away from the other,
so $\Delta b$ and $\Delta e$ have opposite signs.  Off the path, both anchors lie on the
same side of the edge, so distances to both change in the same direction.
\end{proof}

\begin{corollary}[Monotone direction score]\label{cor:mds}
$\mathrm{mds}(v) = |\{u \in N(v) : \Delta b \cdot \Delta e > 0\}|$ equals the number of
incident edges not on the $L$--$R$ path.  Branch body nodes have $\mathrm{mds}(v) = \deg(v)$;
non-junction on-path nodes have $\mathrm{mds}(v) = 0$; junction nodes have intermediate scores.
\end{corollary}

\paragraph{Benchmark results.}
On 192 synthetic path-with-branch graphs, $\mathrm{mds}$ achieves:
AUROC $= 0.960$ on tree instances (no cycles), AUROC $= 0.806$ against
the extended target (branch body and attachment node combined).
On shortcut-containing (cyclic) instances AUROC falls to $0.595$, consistent
with the theorem's tree precondition.
An anchor-free variant using eccentricity-derived anchors (double BFS)
achieves AUROC $= 0.906$ on tree instances (5.4\% gap vs.\ anchored).

\paragraph{Real-graph validation.}
The 2002 NCAA Division~I-A football schedule (115 teams, 613 games) provides
ground truth: the four historically conference-independent programmes
(Notre Dame, Army, Air Force, Navy) are externally identifiable without labels.
Computing $\mathrm{mds}$ on the BFS spanning tree places all four in the top-8
scoring nodes (scores 7--8) with no label supervision, while conference-embedded
teams (e.g.\ Clemson) score zero.

% ====================================================================
\section{Planetary Mean-Motion Resonances}\label{sec:planetary}

\paragraph{Claim.}
Among catalogued mean-motion resonances in planetary systems,
period ratios $p$:$q$ where both $p$ and $q$ are Fibonacci numbers
are observed significantly more often than expected from a uniform or
physically-motivated null model.

\paragraph{Observer projection.}
A resonance pair with period ratio $p$:$q$ (with $p > q$, $\gcd(p,q)=1$)
is classified as \emph{Fibonacci} if both $p$ and $q$ are Fibonacci numbers
(i.e.\ members of $\{1,2,3,5,8,13,21,\ldots\}$), and \emph{non-Fibonacci} otherwise.
This binary classification is the observer projection: it maps the continuous
period ratio to a discrete label using only number-theoretic structure.
The QA mechanism predicts that the Fibonacci shift $T$ (multiplication by
$\varphi^2$ modulo $m$) generates resonance widths
(measured by Laplace coefficients) that are maximised at Fibonacci-ratio
commensurabilities, producing deeper attractor basins and hence higher
empirical frequency.  No resonance frequency data enters this prediction.

\paragraph{Primary catalogue (8 systems, 60 resonances).}
The primary dataset comprises 60 catalogued mean-motion resonances across
8 multi-planet systems with known resonance chains (K2-138, Kepler-80,
Kepler-223, TOI-178, GJ-876, HD~158259, and two solar-system subsystems).
Order-stratified results:

\begin{center}
\small
\begin{tabular}{cccccl}
\toprule
\textbf{Order} & $n_\text{Fib}$ & $n_\text{tot}$ & \% Fib & Exp.\ \% & \textbf{Ratios observed} \\
\midrule
0 & 1 & 1  & 100 & 100 & 1:1 \\
1 & 33 & 43 & 77  & 22  & 2:1, 3:2, \textbf{4:3} \\
2 & 6  & 6  & 100 & 50  & 3:1, 5:3 \\
3 & 4  & 6  & 67  & 40  & 5:2, 8:5, \textbf{4:1}, \textbf{7:4} \\
4 & 1  & 2  & 50  & 33  & 5:1, \textbf{7:3} \\
5 & 0  & 2  & 0   & 25  & \textbf{7:2}, \textbf{9:4} \\
\bottomrule
\end{tabular}
\end{center}

All five statistical tests reject the uniform null (expected Fibonacci fraction
22\% for first-order, 31\% overall):

\begin{center}
\small
\begin{tabular}{lcc}
\toprule
\textbf{Test} & \textbf{Observed} & $p$-value \\
\midrule
All instances (binomial)  & 45/60 (75\%)  & $< 10^{-6}$ \\
First-order only          & 33/43 (77\%)  & $< 10^{-6}$ \\
Exoplanet-only            & 24/31 (77\%)  & $< 10^{-6}$ \\
Fisher combined           & $\chi^2 = 59.0$ & $< 10^{-6}$ \\
Unique ratio types        & 8/14 (57\%)   & 0.040 \\
\bottomrule
\end{tabular}
\end{center}

The most conservative test (unique ratio types, which counts each ratio once
regardless of how often it appears) gives $p = 0.040$.

\paragraph{Robustness to null model choice.}
The uniform null (22\% expected) may under-represent low-integer ratios.
Five alternative null models spanning the plausible range of physical weighting:

\begin{center}
\small
\begin{tabular}{lcc}
\toprule
\textbf{Null model} & Exp.\ Fib.\ frac.\ & $p$ (first-order) \\
\midrule
Uniform ($p \leq 10$, equal weight)          & 22\%  & $4.7 \times 10^{-14}$ \\
Laplace coeff.\ $\alpha^p$, $\alpha=(q/p)^{2/3}$ & 20\% & $1.6\times10^{-15}$ \\
$1/p$                                        & 43\%  & $8.0\times10^{-6}$  \\
$1/(p+q)$                                   & 47\%  & $7.0\times10^{-5}$  \\
$\exp[-(p+q)/10]$                            & 40\%  & $7.5\times10^{-7}$  \\
$1/p^2$ (aggressively low-integer)           & 66\%  & 0.083 (n.s.)         \\
\bottomrule
\end{tabular}
\end{center}

The physics-motivated Laplace-coefficient weighting (derived from the
disturbing-function expansion for first-order resonances under Kepler's
third law; \cite{murray1999,peale1976}) makes the result \emph{more}
significant than the uniform null ($p = 1.6\times10^{-15}$ vs.\ $4.7\times10^{-14}$),
because it further down-weights high-$p$ non-Fibonacci ratios.
Only the aggressively low-integer $1/p^2$ weighting renders the effect
non-significant, and that weighting has no justification from standard
perturbation theory (libration width is approximately constant
across the ratio ladder at fixed eccentricity for first-order resonances).

\paragraph{Out-of-sample replication (341 systems).}
An independent replication using 341 additional planetary systems from the
NASA Exoplanet Archive (raw period data, no resonance labels used in model
building) finds order-1 Fibonacci rate $81.3\%$ ($n = 139$ first-order resonances,
$p = 2.6\times10^{-49}$ vs.\ uniform null).
Fisher's method combining original and replication results across
resonance orders gives $\chi^2 = 288.1$, $p = 1.5\times10^{-54}$.
The replication also shows order-dependence consistent with the QA mechanism:
orders 1--3 remain significantly Fibonacci-enriched ($p < 0.007$);
orders 4--5 are not ($p > 0.7$), matching the predicted fall-off of
resonance width with order~\cite{dale2026fib}.

% ====================================================================
\section{EEG Seizure Prediction}\label{sec:eeg}

\paragraph{Claim.}
QA 4-tuple features derived from topographic EEG clustering add $\Delta R^2 > 0$
to seizure-period prediction beyond a delta-band spectral baseline,
across every patient tested.

\paragraph{Observer projection.}
Ten-second EEG windows from 23 channels are mapped to QA state as follows.
(1)~Compute per-channel Welch power spectrograms.
(2)~Apply $k$-means ($k=9$) to the channel$\times$frequency matrix to obtain
9 spatial cluster centroids.
(3)~Project each centroid to its two dominant PCA components $(c_1, c_2)$.
(4)~Map to QA state: $b = (\lfloor |c_1| \rfloor \bmod 9)+1$,
$e = (\lfloor |c_2| \rfloor \bmod 9)+1$.
(5)~Compute the 4-tuple $(b,e,d,a)$ where
$d = ((b+e-1)\bmod 9)+1$ and $a = ((b+2e-1)\bmod 9)+1$.
(6)~Average the 4-tuple across all 9 centroids.
The result is a 4-dimensional feature vector per window with no
free parameters other than $k=9$ and the PCA truncation,
both fixed before examining seizure labels.
The delta-band baseline (mean delta-band power per channel) is also fixed
before examining labels.

\paragraph{Results.}
Dataset: CHB-MIT Scalp EEG Database~\cite{chbmit}, 10 patients
(chb01--chb07, chb11, chb18, chb20); 350 seizure windows, 412 baseline windows
(balanced by undersampling).
Outcome: McFadden pseudo-$R^2$ of logistic regression on delta-band + QA features
minus logistic regression on delta-band alone ($\Delta R^2$).

\begin{center}
\small
\begin{tabular}{lrrrl}
\toprule
\textbf{Patient} & $\Delta R^2$ & SE & $p$ (LR test) & \\
\midrule
chb01 & $+0.200$ & --- & $1.9\times10^{-5}$ & * \\
chb02 & $+0.177$ & --- & $1.1\times10^{-2}$ & * \\
chb03 & $+0.228$ & --- & $6.6\times10^{-7}$ & * \\
chb04 & $+0.099$ & --- & $2.3\times10^{-2}$ & * \\
chb05 & $+0.277$ & --- & $2.0\times10^{-8}$ & * \\
chb06 & $+0.194$ & --- & $1.1\times10^{-3}$ & * \\
chb07 & $+0.192$ & --- & $1.0\times10^{-3}$ & * \\
chb11 & $+0.293$ & --- & $7.1\times10^{-8}$ & * \\
chb18 & $+0.048$ & --- & $0.107$ & n.s. \\
chb20 & $+0.396$ & --- & $1.4\times10^{-9}$ & * \\
\midrule
\textbf{Mean} & $\mathbf{+0.210}$ & $0.029$ &
  \multicolumn{2}{l}{Fisher $\chi^2=208.0$, $p=2.9\times10^{-33}$} \\
\bottomrule
\end{tabular}
\end{center}

All 10 patients show positive $\Delta R^2$; 9/10 are individually significant
($p < 0.05$, likelihood-ratio test).  The sole exception (chb18, $\Delta R^2 = +0.048$)
still gains.  The median $\Delta R^2 = 0.197$ is close to the mean,
indicating no single patient drives the aggregate.

\paragraph{Feature ablation.}
A companion HI~2.0 classifier experiment (F1-based, separate analysis)
ablates the 10-feature QA set (harmonic-index scalars,
Pythagorean elements $C,F,G,\gcd$, and raw coordinates $b,e$)
across two pre-defined patient cohorts:

\begin{center}
\small
\begin{tabular}{lrrrr}
\toprule
& \multicolumn{2}{c}{\textbf{Stable-positive}} & \multicolumn{2}{c}{\textbf{Hard-negative}} \\
\cmidrule(lr){2-3}\cmidrule(lr){4-5}
\textbf{Feature set} & F1 & $\Delta$F1 & F1 & $\Delta$F1 \\
\midrule
Full (all 10)               & 0.730 & ---      & 0.644 & --- \\
No geometry ($-C,F,G,\gcd$) & 0.693 & $-0.036$ & 0.699 & $+0.056$ \\
No coords ($-b,e$)          & 0.718 & $-0.011$ & 0.644 & $\pm0.000$ \\
Coords only ($b,e$)         & 0.709 & $-0.021$ & 0.740 & $+0.096$ \\
Geometry only               & 0.715 & $-0.015$ & 0.599 & $-0.045$ \\
\bottomrule
\end{tabular}
\end{center}

For stable-positive patients, removing Pythagorean geometry costs
$\Delta\text{F1} = -0.036$; for hard-negative patients the geometry
features actively hurt ($+0.056$ when removed) and raw coordinates
alone outperform the full set by $+0.096$.
The two cohorts differ not only in whether QA features help, but in
\emph{which algebraic layer} carries the signal.

\paragraph{Robustness notes.}
Patients chb08--chb10, chb12--chb17, chb19, and chb21--chb23 (13 patients)
are untouched and available for pre-registered replication.
The observer projection has two unfixed choices ($k=9$ centroids, random seed);
a replication study should fix these before data access and test sensitivity.

% ====================================================================
\section{Climate Regime Classification}\label{sec:climate}

\paragraph{Claim.}
ENSO phase (El Ni\~no / La Ni\~na / Neutral) maps preferentially to distinct QA orbit
families when a five-channel climate index is discretised via the QA observer projection.

\paragraph{Observer projection.}
Monthly values of five NOAA climate indices (ONI, NAO, AO, PDO, AMO; 1950--2023,
$n = 817$ months) are standardised and partitioned into train/test halves.
$K=4$ clusters are fit to the training set via $k$-means; each monthly observation
receives a cluster label.  The cluster centroid is mapped to QA state by taking
$(b,e)$ from the two coordinates that explain the most variance and reducing
modulo 9 ($b,e \in \{1,\dots,9\}$).  The orbit family (Cosmos, Satellite, or
Singularity) is read from the resulting pair using the exact $m=9$ orbit decomposition.
ENSO phase labels (El Ni\~no / La Ni\~na / Neutral) are determined independently
from the Oceanic Ni\~no Index (ONI threshold $|\text{ONI}| \geq 0.5$\,K);
they are not used in the clustering or orbit assignment.

\paragraph{Results.}
The orbit-by-ENSO-phase contingency table ($n=408$ test months) is:

\begin{center}
\small
\begin{tabular}{lrrrl}
\toprule
\textbf{ENSO phase} & \textbf{Singularity} & \textbf{Satellite} & \textbf{Cosmos} & $n$ \\
\midrule
El Ni\~no  & 23.2\% & 37.4\% & 39.4\% & 99 \\
La Ni\~na  &  3.3\% & 96.7\% &  0.0\% & 121 \\
Neutral    & 11.2\% & 74.5\% & 14.4\% & 188 \\
\bottomrule
\end{tabular}
\end{center}

La Ni\~na months concentrate in the Satellite orbit at 96.7\% (117/121 months);
El Ni\~no distributes across all three families with no dominant orbit.
The overall chi-squared test of independence ($\chi^2 = 94.88$, $p = 1.2 \times 10^{-19}$,
$\text{df} = 4$) rejects the null of orbit-phase independence decisively.

A secondary test examines whether the QA Coherence Index (QCI) --- which measures
how consistently consecutive orbit transitions follow the Fibonacci shift prediction ---
partially correlates with future cross-index dispersion after conditioning on lagged ONI.
This partial correlation is $r = -0.135$, $p = 0.007$.
The primary test (orbit distribution) and secondary test (QCI forecasting) are
separate hypotheses; the primary is the load-bearing claim.

% ====================================================================
\section{Cross-Domain Synthesis}\label{sec:synthesis}

\subsection{Summary across domains}

Table~\ref{tab:domains} collects the four applications.
The left columns describe what changes between domains (the observer projection
and the outcome metric); the right columns describe what stays fixed
(the algebraic computation and the modulus).

\begin{table}[h]
\centering
\small
\begin{tabular}{p{2cm}p{3.2cm}p{2.5cm}p{2cm}l}
\toprule
\textbf{Domain} & \textbf{Observer projection} & \textbf{Primary result} & \textbf{Effect} & \textbf{Replicated?} \\
\midrule
Graph networks &
  BFS distances to two anchors $\to$ $(\Delta b, \Delta e)$ sign product &
  AGS theorem (exact) &
  AUROC $0.960$ &
  Theorem \\[4pt]
Planetary resonances &
  Period ratio $p$:$q$ $\to$ Fibonacci label (both integers Fibonacci?) &
  Binomial + Fisher combined &
  $p = 1.5\!\times\!10^{-54}$ &
  Yes ($n=341$) \\[4pt]
EEG seizure &
  23-ch PSD $\to$ $k$-means ($k=9$) $\to$ PCA $\to$ mod-9 $\to$ $(b,e,d,a)$ &
  $\Delta R^2$ vs delta-band baseline &
  $+0.210$ (SE $0.029$) &
  No \\[4pt]
Climate (ENSO) &
  5 climate indices $\to$ $k$-means ($K=4$) $\to$ mod-9 $\to$ orbit family &
  $\chi^2$ test of orbit $\times$ phase independence &
  $p = 1.2\!\times\!10^{-19}$ &
  No \\
\bottomrule
\end{tabular}
\caption{The four domains.  The algebraic layer (Fibonacci shift $T$ at $m=9$,
orbit classification into Cosmos/Satellite/Singularity) is identical across rows.
Only the observer projection and outcome metric change.}
\label{tab:domains}
\end{table}

\subsection{What is literally shared}

The following are identical across all four applications, with no domain-specific
adjustment:

\begin{enumerate}[label=(\arabic*)]
  \item \textbf{The map $T$}: $T(b,e) = (e,\, ((b+e-1)\bmod 9)+1)$.
  \item \textbf{The orbit partition}: Cosmos (72 states, period 24),
        Satellite (8 states, period 8), Singularity (the fixed point $(9,9)$).
  \item \textbf{The modulus $m = 9$} for orbit classification
        and $m = 24$ for the derived coordinates $d$ and $a$.
  \item \textbf{The sign product} $\Delta b \cdot \Delta e \in \{-1,+1\}$
        as the graph-theoretic discriminant (a direct read of the QA step direction).
\end{enumerate}

What differs is the \emph{observer projection}: the function that maps
domain data (period ratios, voltage amplitudes, SST anomalies, graph distances)
into a pair $(b,e) \in \{1,\dots,9\}^2$.  Each observer projection is
defined once, before examining outcome labels, and is not adjusted
post-hoc to improve fit.

\subsection{Why $m = 9$ and $m = 24$ are not free parameters}

Both moduli are constrained by properties that predate any of the four
empirical domains:

\paragraph{Pisano period.}
The Fibonacci sequence modulo 9 is periodic with period $\pi(9) = 24$:
the pair $(F_k \bmod 9,\, F_{k+1} \bmod 9)$ returns to $(1,1)$ after exactly
24 steps~\cite{wall1960}.  No other single-digit modulus $m \leq 9$ has a
shorter Pisano period relative to the number of distinct states.
At $m=9$, the 72-state Cosmos orbit has period exactly $\pi(9) = 24$,
the Satellite has period $8 \mid 24$, and the Singularity is fixed.
The period-24 structure arises from the arithmetic of $\mathbb{Z}[\varphi]/9\mathbb{Z}[\varphi]$,
not from fitting to any domain.

\paragraph{Pythagorean divisibility.}
For every primitive Pythagorean triple $(C,F,G)$ with $C = 2de$,
$F = d^2-e^2$, $G = d^2+e^2$, one has $24 \mid 2CF$~\cite{carmichael1910}.
This is a theorem of elementary number theory: one of $d,e$ is even,
one of $d-e, d+e$ is divisible by 3, and one of $d,e,d-e,d+e$ is
divisible by 4, giving $2CF = 4de(d-e)(d+e)$ divisible by
$4 \cdot 2 \cdot 3 = 24$.  The modulus 24 is therefore the
\emph{natural quantum of the Pythagorean step}, independent of any
empirical application.

Neither derivation involves planetary mechanics, neural signals, climate,
or graph topology.  The fact that $m=24$ is also the length of the
ENSO climatological cycle and the number of standard EEG channels is
a coincidence that the alignment-artifact hypothesis must explain.

\subsection{Alignment artifact audit}

For each domain we identify the specific modelling choice that a skeptic
would point to as a potential alignment artifact:

\paragraph{Graph networks.}
No free parameters exist: the AGS theorem holds for \emph{any} choice of
anchors $L,R$ in \emph{any} finite unweighted tree.  The AUROC $= 0.960$
result on synthetic benchmarks is a consequence of the theorem, not a
tuning outcome.  This domain is \textbf{artifact-free by proof}.

\paragraph{Planetary resonances.}
The Fibonacci label (``both $p$ and $q$ are Fibonacci numbers'') is a
binary predicate with no free threshold.  The uniform null rate
(22\% for first-order) is determined by counting coprime ratios with $p \leq 10$,
a standard convention independent of the data.
The tougher null models (Section~\ref{sec:planetary}) show the result
is robust up to the $1/p^2$ weighting; only a physically unjustified
penalty renders it non-significant.
\textbf{Artifact risk: low.}
The out-of-sample replication ($n=341$ systems, $p = 2.6\times10^{-49}$)
provides independent confirmation.

\paragraph{EEG.}
The choice $k=9$ clusters aligns with the modulus $m=9$, which is not
a coincidence: it was set to match the modulus.  This is a \emph{design
choice}, not a post-hoc tuning, but it means the EEG result cannot be
used to independently validate $m=9$ as the correct modulus.
The result validates that, \emph{given} $k=9$, the 4-tuple features
carry information beyond the spectral baseline.
Separately, all 10 $\Delta R^2$ values are positive; the probability of
this occurring by chance is $(0.5)^{10} \approx 0.001$ under a sign-randomisation
null, suggesting the gain is systematic.
\textbf{Artifact risk: moderate.}  $k=9$ is a confounded choice; held-out
replication with $k \in \{7, 9, 11\}$ would bracket the sensitivity.

\paragraph{Climate.}
The $K=4$ cluster count has no alignment argument with $m=9$ or $m=24$.
The orbit assignment uses a mod-9 reduction of the centroid coordinates,
which does align with the modulus.
The 24-month ENSO cycle does align with $m=24$, but the QCI window in
the secondary test (24 months) was set before the primary chi-squared result
was examined.
\textbf{Artifact risk: moderate.}  The primary chi-squared result
($p=1.2\times10^{-19}$) depends on the centroid-to-orbit assignment,
which is the modulus-aligned step.  A version using $m=7$ or $m=11$
as a placebo would test whether the effect is modulus-specific.

\subsection{Replication scorecard}

\begin{center}
\small
\begin{tabular}{lll}
\toprule
\textbf{Domain} & \textbf{Current status} & \textbf{What would close the loop} \\
\midrule
Graph (AGS)          & Proven (theorem)       & Extension to bounded-treewidth graphs \\
Planetary            & Out-of-sample replicated ($n=341$) & Prospective catalogue update \\
EEG                  & 10 patients, all positive & Pre-registered on held-out 13 patients \\
Climate              & Single dataset, 1950--2023 & Placebo moduli ($m=7, 11$); ERA5 hold-out \\
\bottomrule
\end{tabular}
\end{center}

The joint prior probability that all four results arise by chance (treating
domains as independent) is the product of the individual probabilities,
which is negligible.  However, the four domains are not independent tests
of the same hypothesis in the strict statistical sense: they share the
same analyst, the same algebraic framework, and partly overlapping
design choices.  The appropriate weight to place on the cross-domain
pattern therefore lies between ``four independent replications''
and ``one result tested four ways.''
The planetary out-of-sample replication and the AGS theorem are the
two load-bearing results that are least susceptible to this criticism.

% ====================================================================
\section{Open Questions}\label{sec:open}

\paragraph{Pre-registration and replication.}
The EEG Observer 3 result is the most impressive effect size ($\Delta R^2 = +0.21$)
but also the most complex observer projection.  A pre-registered replication
on held-out CHB-MIT patients (patients 11--23 were not used) with a fixed
analysis plan would substantially increase the credibility of this result.

\paragraph{Orbit theory for higher moduli.}
The orbit decomposition at $m=24$ is empirically known but not fully classified.
A closed-form description of the orbit structure for general $m$ (depending on
the splitting of primes in $\mathbb{Z}[\varphi]$) would explain why $m=24$
in particular generates stable attractor basins.

\paragraph{Scope of the AGS theorem.}
The theorem holds for trees.  Extending it to graphs with bounded treewidth,
or identifying a modified score for cyclic graphs that preserves the
exact separation property, is an open combinatorial problem.
The Koenig gap $2CF$ achieves AUROC $= 0.618$ on shortcut graphs in our benchmark
but does not have a separation theorem behind it.

\bigskip
\noindent\textbf{Acknowledgements.}
The author thanks Ben Iverson and Dale Pond for foundational QA work,
and Norman Wildberger for rational trigonometry and chromogeometry.

% ====================================================================
\begin{thebibliography}{20}

\bibitem{dale2026fib}
W.~Dale.
\newblock Fibonacci preference in mean-motion resonances: evidence from solar and exoplanetary systems.
\newblock Preprint, 2026.

\bibitem{barning1963}
F.~J.~M. Barning.
\newblock Over pythagorese en bijna-pythagorese driehoeken en een generatieproces met behulp van unimodulaire matrices.
\newblock \emph{Math.\ Centrum Amsterdam}, 1963.

\bibitem{carmichael1910}
R.~D. Carmichael.
\newblock Note on a new number theory function.
\newblock \emph{Bull.\ Amer.\ Math.\ Soc.}, 16(5):232--238, 1910.

\bibitem{wall1960}
D.~D. Wall.
\newblock Fibonacci series modulo $m$.
\newblock \emph{Amer.\ Math.\ Monthly}, 67(6):525--532, 1960.

\bibitem{wildberger2005}
N.~J. Wildberger.
\newblock \emph{Divine Proportions: Rational Trigonometry to Universal Geometry}.
\newblock Wild Egg, Sydney, 2005.

\bibitem{chbmit}
A.~L. Goldberger et al.
\newblock PhysioBank, PhysioToolkit, and PhysioNet: components of a new research resource for complex physiologic signals.
\newblock \emph{Circulation}, 101(23):e215--e220, 2000.

\bibitem{murray1999}
C.~D. Murray and S.~F. Dermott.
\newblock \emph{Solar System Dynamics}.
\newblock Cambridge University Press, 1999.

\bibitem{peale1976}
S.~J. Peale.
\newblock Orbital resonances in the solar system.
\newblock \emph{Annu.\ Rev.\ Astron.\ Astrophys.}, 14:215--246, 1976.

\end{thebibliography}

\end{document}
